\section{Related Work --- draft v1}
\label{sec:related}

Five clusters position this work; the two concessions come first,
because a reviewer should meet them in our words before theirs.

\subsection{Structured diagnostics in production compilers}

Production compilers have been converging on diagnostics-as-data by
retrofit. Rust formalized its error format and structured
suggestions in RFC 1644 \cite{rust-rfc1644}, emits machine-consumable JSON
diagnostics \cite{rustc-json}, and maintains Fluent-based message infrastructure
with a public error index \cite{rustc-dev-guide}. GHC's migration from \texttt{SDoc} strings
to "errors as structured values" --- with \texttt{diagnosticReason} and
\texttt{diagnosticHints} fields and public error codes since 9.6.1 --- is
documented across its development tracker and implementers' reports
\cite{ghc-structured-errors,welltyped2021ghc}, and the community maintains an error index \cite{haskell-error-index}. **Concession
one:** the five-role core of \S6.2 is therefore not a new invention
but a distillation; our claims are that designing it in from the
first commit is cheap (\S8), that the isolated knowledge base and
root-cause-granularity taxonomy are the load-bearing choices (\S6,
D3--D4), and that the retrofit pain documented by these projects [4,
9] is evidence the pattern matters. To our knowledge, none of these systems
provides the recorded per-diagnostic execution trace of \S6.3. That
knowledge is not passive: a targeted counterexample sweep (recorded
in the reference registry) found only adjacent work --- trace tooling
for \emph{student programs} \cite{suzuki2017tracediff}, and compiler-internal debug dumps
aimed at the compiler's own developers rather than its users --- and
if later work surfaces, this sentence and C2 narrow accordingly.

\subsection{Optimization transparency and provenance}

LLVM's remarks infrastructure \cite{llvm-remarks} is production-scale prior art for
per-transformation records: passes emit typed remark objects with
source locations, serialized to YAML and visualized with dedicated
tools. \textbf{Concession two:} our per-instruction notes are the same
mechanism at teaching scale. The residue we claim: unification of
provenance with the \emph{diagnostic} architecture under one contract
(one struct family, one renderer discipline, one JSON export), its
measured always-on price (12--17\%, \S9.2 --- LLVM's remarks are opt-in
per build), the human-readable in-text form (\texttt{\# line N:} comments in
the emitted assembly, by analogy to \texttt{.loc} directives but for
readers), and the negative result that per-instruction history is
not transitive under dead-code elimination (\S6.4, \S9).

\subsection{Error-message quality and its contested effectiveness}

A half-century of work examines programming error messages; Becker
et al.'s working-group report maps the landscape \cite{becker2019unhelpful}. Eye-tracking
evidence shows developers genuinely read error messages and that
reading difficulty predicts task performance \cite{barik2017developers}; design guidance
for novice-facing messages continues to develop \cite{denny2021designing}. Yet controlled
studies of \emph{enhanced} messages report effects from ineffectual \cite{denny2014enhancing}
to inconclusive \cite{pettit2017enhanced} --- a contested record extending into the LLM era,
where generated explanations show promise in some deployments \cite{taylor2024dcc}
and ineffectiveness in others \cite{santos2024silver}. This contestation shapes our
claim discipline (\S9.3): we measure component \emph{presence} with a
mechanical instrument descended from Marceau et al.'s evaluation
rubric \cite{marceau2011measuring}, and defer all helpfulness claims to a designed-but-
unexecuted study (\S13). Our contribution is orthogonal to wording
quality: an architecture that \emph{retains the evidence} any wording ---
human-authored or LLM-generated \cite{taylor2024dcc} --- would need.

\subsection{Error localization and interrogative debugging}

Constraint-based localization computes \emph{where} the likely mistake
is: general static-error diagnosis via Bayesian-ranked unsatisfiable
constraints \cite{zhang2014general,zhang2017sherrloc}, type-error diagnosis \cite{zhang2015diagnosing}, and SMT-based
localization \cite{pavlinovic2015smt}. These are complementary: they improve one field
of our diagnostic (the span and root-cause selection), while our
architecture governs how the whole explanation is assembled,
preserved, and delivered. The Whyline \cite{ko2008debugging,ko2010extracting} pioneered
why/why-not questions over \emph{program} executions via tracing and
slicing, a family that extends into education through trace-
divergence debugging of student code \cite{suzuki2017tracediff}; we ask the interrogative
question of the \emph{compiler's own run} --- a far shallower trace,
recorded live at negligible cost (\S9.2) rather than reconstructed by
post-hoc analysis.

\subsection{Compiler visualization in education}

The lineage is old --- the UW Illustrated Compiler graphically
displayed its own control and data structures to students in 1988
\cite{andrews1988uw} --- and continues through AST construction animators \cite{almeida2008vast},
interactive stage simulations \cite{rodriguez2013simulations}, and, recently, JIT IR-evolution
explorers \cite{dalbo2025jitscope}. Our visualizer (\S7)
sits in this lineage as a \emph{consumer} of the architecture rather than
a contribution: its distinguishing property is that every panel
renders the same structured values the CLI renders, because the
architecture guarantees a single source of truth (\S6.2, D6).

\subsection{Positioning summary (Table T2)}

\begin{table}[t]
\centering
\small
\begin{tabular}{lllll}
\toprule
Capability & GHC \cite{ghc-structured-errors} & rustc \cite{rust-rfc1644,rustc-dev-guide,rustc-json} & LLVM \cite{llvm-remarks} & this work \\
\midrule
Structured diagnostics as values & retrofit & retrofit & --- & designed-in \\
Root-cause-indexed public error index & \checkmark{} \cite{haskell-error-index} & \checkmark{} & --- & \checkmark{} (taxonomy = index, D3) \\
Machine-readable diagnostic export & \checkmark{} & \checkmark{} \cite{rustc-json} & \checkmark{} (remarks) & \checkmark{} \\
Per-diagnostic \textbf{recorded execution trace} & --- & --- & --- & \checkmark{} (\S6.3) \\
Optimization provenance records & --- & --- & \checkmark{} opt-in \cite{llvm-remarks} & \checkmark{} always-on, unified, priced \\
Diagnostic-quality \textbf{property} tests & snapshot-style & golden \texttt{.stderr} \cite{rustc-dev-guide} & --- & property-level + mutation-measured (\S9) \\
\bottomrule
\end{tabular}
\end{table}



